\subsection{F.1: sources, exact scope, and search outcomes}

\paragraph{Status of the audit.}
The primary statements below were opened at the locations recorded
in the bibliography. Two requested original bodies have not been
obtained: Isbell (1974) and McDuff (1979).
Their publisher or bibliographic records are identified explicitly;
they are not treated as checked proofs or as absent results.
This access gap is why F.1 remains incomplete. None of the new
arguments in this return depends on an unchecked assertion from
those two bodies. Oberst's original 1968 article was subsequently
obtained from the G\"ottingen digitization and checked directly,
including the hypothesis and the conjectural extension discussed below.

\paragraph{Stabilizers and imaginaries.}
Macpherson--Tent, Lemma~2.7(i) of \cite{ABMacphersonTent}, states
$G_{A\cap B}=\langle G_A,G_B\rangle$ for finite $A,B$ in a free
homogeneous relational structure. Part~(ii) gives a unique least
finite support for a finitely supported invariant relation; the
preceding paragraph identifies the resulting weak elimination of
imaginaries in the $\omega$-categorical case. Lemma~2.8 excludes
proper nontrivial open normal subgroups. It does not assert (Op):
the random graph already has open setwise stabilizers of finite
sets, as our graph example verifies. Its introduction specifically
identifies posets and Hrushovski constructions as cases not covered
by its method. Our ordered and strong versions use the finite
bridge proofs in Part B, with the intersection hypothesis made
explicit.

Conant's Definition~2.3 and Theorem~5.6 in \cite{ABConant2017}
concern a theory admitting an independence relation satisfying
invariance, monotonicity, symmetry, full transitivity, full existence,
stationarity, freedom and closure. The conclusion is elimination of
hyperimaginaries and weak elimination of imaginaries. It is not a
blanket assertion of (Op) for arbitrary strong classes with an added
linear order. We checked the current arXiv version, including its
2023 corrigenda; the corrections concern later material, not a
replacement for Theorem~5.6. The relevant finite-support argument
was also read directly.

Evans--Ghadernezhad--Tent, Theorem~4.15 of
\cite{ABEvansGhadernezhadTent2016}, proves simplicity of
$\Aut(M_0/\operatorname{cl}^d(\varnothing))$ for their uncollapsed
predimension $\delta=n|V|-m|R|$, with $m,n$ positive coprime,
when $r=2,n>m$, or $r\ge3,n\ge m$; each nonidentity element normally
generates in at most $96$ conjugates or inverse conjugates.
Theorem~5.10 gives the $192$-conjugate conclusion for the stated
monodimensional $M_f$ under Assumption~5.5. Neither statement is
(St), (Op), or a general weak-imaginary-elimination theorem for our
ordered positive-predimension classes. The text was searched for
stabilizers, imaginaries, intersections and normal subgroups, and
the relevant closure and independence sections were read. Section~2
explicitly uses a closure topology that need not equal the usual
finite-point-stabilizer topology; transferring its topological
arguments without checking this would be invalid.

\paragraph{Amenability, Ramsey expansions, and EPPA.}
Kechris--Rosendal, Proposition~6.4 of \cite{ABKechrisRosendal2007},
identifies the Hrushovski property (EPPA) of an ordinary Fra\"iss\'e
class with the existence of increasing compact subgroups having
dense union in its automorphism group. Both directions of its
proof were read. The strong-class implication is also stated as
Theorem~2.24 in \cite{ABEHN2019}. Our averaging proof in
Proposition~\ref{A:averaging} derives rational exactness from precisely
that compact-chain hypothesis.

Angel--Kechris--Lyons, Theorem~1.4 of \cite{ABAngelKechrisLyons2014},
proves unique ergodicity for the automorphism groups of a countably
infinite-dimensional vector space over a finite field, the random
$\mathcal A$-free hypergraphs in their finite hypergraph language
($\mathcal A$ consists of irreducible forbidden structures), and
$U_S$ for a countable additive subsemigroup $S\subseteq\R_{>0}$.
Proposition~9.3 gives amenability when all linear orders form a
companion class. Theorem~13.3, for an EPPA class and a reasonable
order expansion, equates unique ergodicity of its expansion flow
with its quantitative ordering properties; for an excellent pair
this is unique ergodicity of the group. Here ``Hrushovski class''
means EPPA, not an arbitrary predimension construction. These
results do not assert amenability of the independent-order
stabilizer of the vector-space group; Theorem~\ref{A:ordered-vector-amenability}
provides that additional proof here.

Bodirsky's Theorem~1.5 and its proof in
\cite{ABBodirsky2014}, pp.~4 and 9--11, require hereditary
relational classes with strong amalgamation. Corollary~1.4(3),
pp.~3--4, applies to our poset expansion. We also inspected the
affine-vector entry in Definition~1.7 and Corollary~1.8, p.~5;
its listed second factors are tournaments, graphs, and
$K_n$-free graphs, not an independent linear order. We do not
use it to assert that natural orders together with arbitrary
reference orders on vector spaces are Ramsey. Nor can the
strong-amalgamation hypothesis of Theorem~1.5 simply be
transferred from closed vector subspaces to the hereditary
relational age: over two independent vectors in characteristic
two, the third vector given by their sum is forced, so two
copies containing that third vector cannot be strongly
amalgamated over just the pair. The Ramsey statement for natural
orders together with arbitrary reference orders remains unproved
here, but it is no longer needed for rational exactness of this
example. Theorem~\ref{A:ordered-basis-ramsey} instead derives a
Ramsey category of distinguished bases with independently ordered
nonempty subsets from Hubi\v cka--Ne\v set\v ril,
\cite[Theorem~4.46, arXiv v4, pp.~75--76]{DHubickaNesetril2019}.
The opened statement and incidence-closure proof allow independent
orders on vertices and on relation tuples. Our padded-tuple completion
checks the passage to all subset sums, and the exact-category
finite averaging in Theorem~\ref{A:ordered-vector-rational-exactness}
then proves rational exactness.

Kechris--Soki\'c, Proposition~2.1 and Remark~3.2 of
\cite{ABKechrisSokic2012}, supply the poset dynamical obstruction
used for comparison in Part A. Our four-point proof treats the
independent reference order explicitly. Moore's published
Theorem~7.1(5)--(6) in \cite{ABMoore2013}, numbered 6.1 in the
supplied preprint, is the convex Ramsey--amenability equivalence;
our proof extends it to the precise strong setting. These positive
probability statements cannot replace the signed affine equalizers
that characterize algebraic exactness.

Zucker's requested 2016 paper is \cite{ABZucker2016}:
Theorem~8.14 equates metrizability of the universal minimal flow,
finite Ramsey degrees of every finite structure, and a reasonable
precompact expansion having the Ramsey and expansion properties.
Corollary~8.15 gives the generic point property. This yields
Corollary~\ref{A:metrizable-umf-exact}, but does not derive
metrizability from amenability. The separate 2014 paper
\cite{ABZucker2014}, Theorem~3.1, gives a finite-expansion positivity
criterion: every nonzero vector in the rational span of the
restriction-relation vectors has coefficients of both signs if
and only if the group is amenable. Its hypotheses include an
excellent pair. It is not a theorem that an arbitrary algebraic
invariant functional is positive.

\paragraph{Toposes and smooth modules.}
Caramello's Theorem~2.9 in \cite{ABCaramello2008} applies to a
consistent theory of presheaf type whose finitely presentable
models have AP, JEP, and a finite or countable dominating family.
Under its stated monomorphism, countability and infinite-model
hypotheses it produces the homogeneous countable model and
\[
 \Set[T']\simeq
 \operatorname{Sh}((\mathrm{f.p.}\,T\text{-}\mathrm{mod})^\op,J_{at})
 \simeq\operatorname{Cont}(\Aut(M)).
\]
The displayed equivalence is of atomic sheaf toposes. It is not
an equivalence between all covariant diagram categories, and hence
alone does not prove the unqualified group invariance of totality.
The proof and the subsequent Galois-property proposition were read;
Theorem~\ref{A:totality-finite-closure} supplies a separate verified
transfer under finite strong closure.

Tsankov's Theorem~1.3 in \cite{ABTsankov2012} classifies continuous
unitary representations of an oligomorphic group: every irreducible
is induced from an irreducible representation of a finite quotient
$C(V)/V$, where $V$ is open and normal in its commensurator, and
every unitary representation is a sum of irreducibles. This Hilbert
space theorem is not exactness of arbitrary smooth algebraic
coinvariants. The latter dictionary is Theorem~\ref{A:coinvariants}.
Li, \cite{ABLi2025}, Theorem~3.6, proves that restriction along a
dense subgroup gives an equivalence of discrete representation
categories. Theorem~4.4 identifies the opposite canonical age with
the point-stabilizer orbit category under strong amalgamation.
We retain that hypothesis: for a perfect matching, a point and its
matched pair have the same pointwise stabilizer, so there is a
reverse orbit map without a corresponding pair-to-point embedding.
Thus the introductory assertion of full faithfulness without this
qualification cannot be used. Theorem~1.5 concerns highly
homogeneous actions and a polynomial coefficient ring; it supplies
no general Noetherian theorem for the groups considered here.
The abstract of \cite{ABCastellanoWeigel} concerns rational discrete
\emph{cohomology} of totally disconnected locally compact groups;
its full proof was not imported, and neither local compactness nor
the change from cohomology to coinvariants is automatic here.

The current primary version of Evans, \cite{ABEvans2026}, is
arXiv v2 of 9 September 2026. Lemma~3.7(i), including its proof,
states that invariants preserve exact sequences of compact
continuous $G$-modules when $G$ is extremely amenable; a compact
coset fibre has a fixed point. Theorem~1.4 and Corollary~1.5 give
finite-stabilizer orbit-augmentation tests and definable separation
for $\omega$-categorical Ramsey structures. Their Ramsey hypothesis
cannot be replaced here by amenability without a new argument.

Yang--Boja\'nczyk--Klin, \cite{ABYangBojanczykKlin2026},
Theorem~5.4 and Corollary~5.5, prove finite length in characteristic
zero under oligomorphic approximation, including unordered vector
atoms. Their Theorem~7.3 treats free amalgamation in a finite
unary/binary relational language, with an optional generic order.
Neither statement applies directly to the independently ordered
vector-space class: its finite ordered structures are rigid, so
they cannot have the uniformly bounded automorphism-orbit counts
of Definition~5.1, and linear span is a nontrivial algebraic closure.
Moreover, finite length alone is not used here to assert boundedness
of an invariant functional on a permutation submodule. The
introduction explicitly leaves general equivariant Noetherianity
unsettled; the conditional saturation hypothesis in
Proposition~\ref{C4:noetherian} therefore remains a hypothesis.

\paragraph{Exact colimits and the torsor formulation.}
The supplied primary Isbell--Mitchell papers were read directly.
Theorem~1 and Lemma~1 of \cite{ABIsbellMitchell1973}, and
Theorem~3.1, Lemma~3.2 and Proposition~3.3 of
\cite{IsbellMitchell1976}, say that integral exactness is equivalent
to filteredness of the components of the affine envelope, or to
condition I and affine equalizers of every finite parallel family.
The explicit remark following Proposition~3.3 extends this to
commutative coefficient rings. Their Theorem~4.1 is the fixed-point
statement for endomorphisms of set actions satisfying their
condition IV; it is not the compact abelian torsor assertion of
Lead~3. Theorems~\ref{C:duality} and \ref{C:torsors} prove that
assertion here by bar--cobar duality and the affine cocycle formula.

Oberst, \cite{ABOberst1968}, was read in the original printed
pages~87--115 from the G\"ottingen digitization. Theorem~I.2.3,
pp.~93--94, is a \emph{homological} cofinality criterion: the
cofibres $y\backslash X$ are connected and integrally acyclic if
and only if restriction induces all derived-colimit isomorphisms,
for every abelian coefficient category with exact direct sums;
testing $\Ab$ suffices. It does not replace acyclicity by weak
contractibility. Theorem~II.2.1, pp.~109--110, identifies integral
exactness with exactness for every AB5 abelian coefficient category
and with filteredness of the indicated element category over the
additive envelope $\overline{\Z X}$ and constant $\Z$.

Definition~II.1.2, p.~108, gives (F3), which is our condition I,
and (F2$'$): a parallel pair that becomes equal after
\emph{precomposition} also becomes equal after some
\emph{postcomposition}. Theorem~II.4.1, pp.~111--114, proves that
integral exactness implies (F3), and that, \emph{under (F2$'$)}, it
is equivalent to filteredness of every connected component.
The proof first treats subgroupoids and cyclic groups, then uses
the calculus of left fractions. Corollary~II.4.3, p.~115, records
that the groupoid completion of any integrally exact category has
simply connected components. These statements were checked
against the proof and the original page images.

Remark~II.4.4, p.~115 (also the introduction, p.~88), conjectures
removing (F2$'$) and says that the arbitrary monoid case would
suffice. The extension is false: the exact, nonfiltered category
$\Delta_{\rm inj}$ is already a counterexample. It fails
(F2$'$), since two distinct injections agreeing at a chosen
endpoint can be preequalized, whereas injective postcomposition
cannot equalize them. The proposed monoid sufficiency is therefore
not a valid reduction to retain: the monoid assertion is true by
the separately checked Isbell note below. No all-prime-to-rational
exactness theorem is asserted in Oberst's article.

Isbell, \cite{ABIsbell1974}, was located at the Annals DOI record
and JSTOR entry, but the original five-page body was not retrieved.
No theorem number from that body is certified by this audit.
The separately opened four-page 1968 note of Isbell,
\cite{ABIsbell1968}, has an unnumbered theorem in Section~1,
pp.~569--570, giving a combinatorial exactness criterion, and an
unnumbered theorem in Section~2, pp.~571--572, proving that a
monoid index has integral exact colimits precisely when it is
filtered. This historical monoid-index theorem is not a result
about free-monoid-valued diagrams, and the note is not a substitute
for the requested 1974 paper.

\paragraph{The orbit-category comparison.}
L\"uck, Definition~1.8 and Theorem~1.9 in \cite{ABLueck2005}, give
the classifying-space characterization by isotropy in a family and
contractibility of all corresponding fixed-point spaces. Section~5.2
sets up the orbit-category chain functor and Bredon invariants.
Nucinkis, Section~2 of \cite{ABNucinkis2004}, defines the tensor
product of left and right orbit modules. Theorem~3.2 states that a
right orbit module is flat exactly when it is a filtered colimit of
finitely generated free orbit modules; equivalently every map to it
from a finitely presented module factors through a finitely generated
free module.
The usual Bredon homological dimension uses the constant
contravariant module. The tensor for our covariant diagrams over
$A\simeq\mathcal O^\op$ uses the constant \emph{left} orbit module.
The variance correction and the required fixed-space check are
proved in Propositions~\ref{B:orbit-dictionary}--\ref{B:orbit-model};
a standard Bredon dimension cannot simply be substituted.
L\"uck's survey assumes local compactness; our open-isotropy
simplicial model is proved directly without that assumption.

\paragraph{Discrete monoids and the free-monoid question (source audit).}
McDuff's 1979 paper \emph{On the classifying spaces of discrete
monoids} is the classical realization result: arbitrary
path-connected weak homotopy types occur as classifying spaces
of discrete monoids \cite{GMcDuff1979}.  In this run the
bibliographic record and reproduced abstract at
\url{https://celebratio.org/McDuff_D/article/412/} were opened,
but the original article body was not obtained.  Its proof and
original theorem numbering are therefore not certified here,
and no argument in this return depends on that proof.  The
realization statement by itself gives no criterion for a
specified colimit of free monoids.

Fiedorowicz's primary 1984 PDF was opened and Proposition~4.4,
including its proof, was checked \cite{GFiedorowicz1984}.
For a discrete monoid $P$ with group completion $P^{\mathrm{gp}}$,
it relates asphericity of $BP$, the comparison with
$BP^{\mathrm{gp}}$, and vanishing of the positive Tor groups
of the group-completion module.  It is a target-asphericity
criterion, not a proof that a particular homotopy-colimit map
is a weak equivalence.  Theorem~4.1, Remark~4.2 and
Lemma~5.2(a) were also read.  The amalgamation argument cannot
be imported without repairing the flatness step: the opened
published paper of Gray--Steinberg gives an explicit
counterexample to that lemma in Remark~4.11
\cite{GGraySteinberg2024}.  Their example uses the amalgam of
the infinite cyclic group with the free monoid on two letters
over the free monoid on one letter.  The publisher's 2025
erratum was opened as well; it corrects submission dates and
does not remove this mathematical warning.  The primary PDF
URLs and exact opened locations are recorded in the cited
bibliographic entries.

Searches included ``free monoids homotopy colimit'', ``exact
colimits free monoids'', ``totally aspherical free monoids'',
``locally free monoids homology'', and the author/title searches
for these three sources.  They did not locate a statement
equivalent to the prompt's Part~G.1, either the specified gap
comparison or the general existence question for nonzero
$L_1\colim M_{\mathrm{ab}}$ on $\Cdelta$.
Nor did they supply an independent published equivalent of the
exact-index comparison DF~D.6.  This is the outcome of the
documented search, not a claim of priority or a declaration
that the unresolved assertions are open in the literature.
The gap comparison is proved in Theorem~\ref{G:gap}; the
general Part~G.1 assertion remains explicitly undecided.


\paragraph{The scope of oligomorphic integration.}
We opened both the full arXiv version and the published text of
Harman--Snowden \cite{DHarmanSnowden2026}. Definitions~10.5--10.6
construct a measure-dependent completed algebra $A_k(G;\mu)$ and
its smooth-module category. Remark~10.10 explicitly distinguishes
this from the category of all smooth $k[G]$-modules: open
stabilizers do not by themselves supply the stronger completed-algebra
smoothness. Theorem~13.2 concerns this additional integration
structure and, under its regularity hypotheses, gives finite-length
and semisimplicity results there. It therefore does not establish
exact coinvariants for the unrestricted smooth modules in Part A,
or the automatic boundedness of rational separating functionals
needed by the remaining amenability argument. We import no such
conclusion from this source.


\paragraph{Search outcomes for the requested leads.}
Author/title and exact-phrase searches included ``smooth
coinvariants exact Fraisse'', ``amenability rational coinvariants'',
``exact colimits compact abelian torsors'', ``Fraisse total
asphericity open subgroups'', ``Bredon flat dimension open
stabilizers'', and the monoid searches recorded above.
For Lead~1, the opened sources supply the related atomic-topos and
unitary-representation statements, but no located statement was
the full all-smooth-module exactness equivalence of
Theorem~\ref{A:coinvariants}. For Lead~3 the standard duality and
cocycle mechanisms were found, but the exact combined formulation
for all small index categories was not located in the requested
papers. For Lead~7 we found the finite-support stabilizer results,
not the complete $E_U$ homotopy reduction with its empty-support
exception. Lead~12 has the established orbit-category framework,
with the variance and fixed-space qualifications just stated.
For A.3, the searched sources give the restricted positive cases
and the familiar amenability obstructions, but did not supply the
signed-poset or total ordered-vector separators proved here, or
settle the remaining implication (iii). The G.1 search outcome is
recorded in the preceding monoid paragraph. These are bounded
search reports, not claims that no equivalent result exists.
